\chapter{Recovering Proofs and Using Negative Results}
\label{ch:proof-practice}

\section{How to take over a research decision}
Start with the mathematical statement in this book. Write down the
carrier, hypotheses, domain, and conclusion you need. The source
declaration and tested revision then determine whether an existing
theorem can supply it. Read the original consumer when the planned
use is stronger than the theorem's name suggests. For a result
carried across repositories, check the literal source and imports
in the intended tree. An unchanged theorem body can still have a
different dependency closure.

The dated atlas indexes the results that enabled the exposition.
The live checkpoint gives the current decision state. In the November
22 snapshot the M4 Lie-defect centrality route and the physical
massless multiplier have distinct receipts. W0C2B1 remains the last
accepted M4 anchor and W0C2B2 is blocked. GM4B2 qualifies the graph
closure on the physical space. The twenty-seven M4 anchors belong
to that register. Chapter~\ref{ch:frontier} locates the current
scientific joins.

\section{Two examples of recovering a usable theorem}
In P2-C4B the hard part was integrability of the actual smooth-core
energy densities. A small equality let the proof use existing
source integrability results. Before reopening a long algebraic
expansion, search for a source lemma whose conclusion already
matches the term after a named rewrite.

In D0b1 the generic graph argument and the paired specialization
failed as a combined proof attempt at the paired join. D0c0 isolated
and qualified the generic partial operator. D0c1 consumed it using
a named source action on paired data. This found a proof that could
be reused on the original carrier while preserving the failed
attempts and their exact meaning. H3B2's use of first projection
from the graph is a related example on the charge domain.

\section{What each failed test tells us}
A parser, import, pin, or cache error reports a failure to evaluate
the intended proof. A failed elaboration reports an attempted
implementation that Lean could not check. A failed independent
consumer can reveal that a general theorem's source binding is
missing or too costly to elaborate. A mathematical counterexample
must exhibit an admissible object and a failed conclusion. These
outcomes lead to different research decisions. The accepted
record should preserve the tested source, first failure, and the
reason the next attempt changes the proof strategy.

Ask what a successful result would enable before staging it.
A candidate should remove a named obstruction in the argument,
test a necessary premise, or provide an exact counterexample.
The most informative next step can be a comparison of existing
declarations, a bounded analytic test, or an independently
compiled source consumer.
